\chapter{第九讲\quad 排除犯罪事由I：排除违法事由}
\label{lecture:09}
\noindent{\small\hyperref[ref:lecture:09]{查阅本讲扩展规则}}\par

先以同一行为、同一相对人及同一受损法益为评价对象。若行为表面上符合某罪的构成要件，才继续问它是否因防卫、避险或承诺而被允许。一把椅子打中侵害人，同时毁损无辜第三人的椅子，要分开写两条评价路径。

\section{第一节\quad 正当防卫}

\paragraph{法条和模型}
《刑法》第20条。令 \(B(d,a,t)\) 表示防卫行为 \(d\) 针对 \(a\) 在 \(t\) 时的现实不法侵害，并具有制止该侵害、保护合法权利的性质。为便于答题，将基础条件分为 \(U\)（现实不法侵害）、\(T\)（正在进行）、\(S\)（防卫指向）、\(O\)（反击对象与侵害相连）。分别验证这四项后才进入限度审查。令 \(M\) 为明显超过必要限度，\(H\) 为该防卫行为造成侵害人重伤或者死亡，\(V\) 为正在进行的、符合第3款的严重暴力侵害。普通防卫和特殊防卫的分支为
\[
\begin{aligned}
B\land V &\Longrightarrow \text{对侵害人的伤亡适用第20条第3款},\\
B\land\neg V\land M\land H &\Longrightarrow \text{防卫过当},\\
B\land\neg V\land\neg(M\land H)&\Longrightarrow \text{一般防卫排除违法}.
\end{aligned}
\]
过当以后，仍须另检验相应罪名的故意、过失等条件。依法构成犯罪时，应当减轻或者免除处罚。若行为本身不具防卫性质，例如故意挑起攻击后借机报复，第一步便不能推出 \(B\)。

\paragraph{推导}
乙持刀追刺甲，甲击打乙持刀的手使刀掉落，乙因这一击骨折。若事实证明当时持续面临刺杀的现实危险，击打是有效而相当的制止方法，则
\[
\begin{aligned}
&\text{持刀连续追刺}\Longrightarrow U\land T,\\
&\text{打落刀以制止追刺}\Longrightarrow S\land O,\\
&U\land T\land S\land O\Longrightarrow B,\\
&\text{行为时可用手段、危险和击打部位}\Longrightarrow \neg M.
\end{aligned}
\]
即使骨折经鉴定达到重伤，\(\neg M\land H\) 仍推出 \(\neg(M\land H)\)。若追刺事实还满足严重危及人身安全的行凶标准，则可进一步适用第3款。刀落地后乙仍可立即捡刀反击，须重新判断 \(T\)。若乙已失去继续侵害的能力，甲事后再打他，第二次击打因 \(\neg T\) 离开防卫分支。

\paragraph{边界和分歧}
《正当防卫指导意见》要求从防卫时的具体处境审查侵害起止。侵害人已拿到财物但还能当场追阻追回，可能仍满足 \(T\)。单凭先动手或双方均受伤，不能把案件机械归入互殴。假想防卫是 \(U\) 为假而行为人误以为为真，转入事实错误与过失审查。偶然防卫是客观制止侵害但行为人没有防卫认识。教材列出两套模型。认识不要说可保留防卫结论。认识必要说不承认防卫，并另审查危害结果的规范归属。教材按后一模型推演的特定故意型案件可以得到未遂，过失型案件可因缺乏可归责的实害结果而不构成过失犯罪。这个结论依赖所用理论和案情。面对复合动机，要分析行为是否实质指向制止侵害以及是否借防卫之名报复。

\section{第二节\quad 紧急避险}

\paragraph{法条和模型}
《刑法》第21条。令 \(D\) 为现实危险正在发生，\(P\) 为行为指向保护合法权益，\(N\) 为当时没有合理可行且损害更小的替代措施，\(L\) 为损害没有超过必要限度。令 \(Q\) 为行为人因职务或业务负有特定责任，\(W\) 为只用避险规则免除自身应承担的相应危险。则
\[
D\land P\land N\land L\land\neg(Q\land W)
\Longrightarrow \text{依法成立紧急避险}.
\]
危险可以源自自然力、动物或人的行为。教材用被保全法益与被牺牲法益的具体性质、强度和情境比较来判断 \(L\)。在可替代的财产利益之间，有时可容许相当损害。生命的价值不能按人数或单价直接相加。若明显有其他合理方法却伤害第三人，先在 \(N\) 处停止。适用第21条第2款须另证超限和不应有的损害，构成犯罪时应当减轻或者免除处罚。

\paragraph{推导}
甲发现孩子被急流困住，唯一能及时拖住孩子的工具是乙锁在岸边的救生绳。现场没有其他绳索，呼救来不及，甲剪开乙的普通锁并救人。
\[
\begin{aligned}
&\text{急流中面临溺亡}\Longrightarrow D,\\
&\text{为救孩子取绳}\Longrightarrow P,\\
&\text{无及时、较小损害的替代方法}\Longrightarrow N,\\
&\text{以普通锁的小额损害避免紧迫的生命危险}\Longrightarrow L.
\end{aligned}
\]
甲也没有落入 \(Q\land W\) 的特定责任例外，故可就毁锁适用避险。若箱门本来未锁，且甲知道直接开门同样及时，\(\neg N\) 会阻止对剪锁的避险推导。

\paragraph{对照}
暴力侵害人挥棍时，甲对其反击，按防卫审查。甲拿第三人的器具抵挡并致器具损坏，第三人的权益受损部分另按避险审查。若侵害人自己把第三人的花瓶扔向甲，甲本能阻挡，先问破损是否应归责于甲。归责不能建立时，无须再为甲寻找正当化事由。两人同时遇险，一人为保全自己而杀死同样有获救机会的无辜者，不能把两条生命代入普通财产等值规则。

\section{第三节\quad 被害人承诺}

\paragraph{教材模型}
刑法总则没有把被害人承诺设为普遍免责条款。令 \(R(r,v)\) 为法益 \(r\) 可由权利人 \(v\) 处分，\(C(v,r)\) 为有理解该处分意义的能力，\(T\) 为同意在实施时有效，\(I\) 为同意真实且没有影响处分内容的重大事实错误，\(F\) 为行为符合对象、方式与风险范围。写为
\[
R(r,v)\land C(v,r)\land T\land I\land F
\Longrightarrow \text{就 }r\text{ 可排除相应的侵害评价}.
\]
若罪状本身以违反意愿为要件，同意首先可能使该要件不成立。若尚存在其他人的权益或公共法益，必须另算。对重伤、生命和医疗处置，要进一步查法益能否处分和专门规范，不能仅凭私人同意推定合法。教材的推定承诺模型要求当时不能及时获得实际意思，并有具体根据推认权利人会同意，已知的明确拒绝优先。

\paragraph{推导}
甲有处分自己废旧桌子的权限，清楚知道乙将拆掉桌面，在拆除前明确同意，乙依约仅拆桌面。由所有权和处分范围得 \(R\land F\)，由真实且及时的具体同意得 \(C\land T\land I\)，从而就该桌面的毁损可排除刑事侵害评价。若乙另拆共同所有的柜门，甲无权独自同意拆柜门，相对于该法益的 \(R\) 或 \(F\) 不成立。若乙骗甲说只卸下可复装的螺丝，实际锯断桌面，错误直接涉及处分内容，不能以此前的表示证明 \(I\land F\)。

\paragraph{事实错误与动机错误}
被欺骗而误判行为对象、内容或实质风险，可能破坏 \(I\)。行为人知道将拆桌面，只因误信以后会收到一封感谢信才表示同意，错误通常只在动机，仍需判断有无条件和范围限制。紧急把昏迷者的电脑移开以防进水，可根据具体信息审查推定承诺。若对方已明确禁止开启加密文件，不能凭猜测推翻该拒绝。

\section*{怎么记与闭卷自测}
先写受损对象，再按三条线复述。\textbf{防卫看因、时、意、对象、限}，限度牢记 \(M\land H\) 同时成立才进入普通防卫过当分支。\textbf{避险看危、护、不得已、衡量、特定义务}，把替代方案写出来。\textbf{承诺看权、能、时、真、范围}，每变动一个对象就重新检查处分权。闭卷先画三条判定线，再用反例把每条线中的一个条件划掉。
\begin{enumerate}
\item 侵害人失去反击能力后，甲另一次报复致其重伤。为何先问 \(T\) 而非 \(M\land H\)？\par
答：第二次行为的时间条件不成立，无法进入 \(B\)，重伤并不能补齐防卫基础。
\item 甲事先同意乙烧自家物品，火势会蔓延到相邻住宅。该同意可以证明什么？\par
答：在同意真实且范围内，可评价自家财产权益的处分。邻居权益和公共安全的审查仍须独立进行。
\end{enumerate}
